\section{The one-sided envelope with both unused products}\label{lc:left}
For \(0\le i\le j<m\), the exact left summand is
\begin{equation}\label{lc:left-summand}
 A_{\rm rest}(u)b^ic^{d-j}
       \left(\prod_{r<i}\alpha_r\right)
       \left(\prod_{j<r<m}\alpha_r\right)K_{i,j}(u),
\end{equation}
where \(A_{\rm rest}=\prod_{r=m}^d\alpha_r\) is the whole omitted
upper-half product. On the real interval it has no positive
exponential cost. At \(N=2m\),
\[
 A_{\rm rest}(u)=2(-1)^mT_m(\sigma/2),\qquad |A_{\rm rest}|\le2.
\]
Its roots and leading coefficient prove the identity.
At \(N=2m+1\), it is, up to sign, \((\sigma+2)P_m(\sigma)\), where
\[
 P_m(2\cos t)=\frac{\cos((m+1/2)t)}{\cos(t/2)}.
\]
The quotient is a finite alternating sum of \(m\) paired exponentials
and one constant, so \(|P_m|\le2m+1=N\).
Hence \(|A_{\rm rest}|\le4N\), including the removable endpoint.
These are the original unused products, not auxiliary normalizations.

For \(b,c>0\), apply the transfer and Riemann estimates separately
to the two outside products in \eqref{lc:left-summand}. Only then
use \eqref{lc:prefix-zero}. With \(\mu=i/N,\nu=j/N\), the upper
action becomes
\begin{equation}\label{lc:left-action}
\begin{split}
 F_{\mathrm L}(\mu,\nu,\delta)
 &=\mu\log b+(1-\nu)\log c+
                   \int_\mu^\nu\log L_\delta(t)\,dt,\\
 L_\delta(t)&=\frac{\lambda_\delta(t)}{|a_\delta(t)|},
                      \qquad0\le\mu\le\nu\le\tfrac12.
\end{split}
\end{equation}
This quotient is only continuum notation. No finite division by
\(\alpha_r\) occurred. Moving an endpoint by \(O(N^{-1})\) changes
the integrals by \(o(1)\) uniformly, by logarithmic integrability.
There are at most \(N^2\) pairs \((i,j)\), contributing only \(o(N)\).
If \(b=0\), the finite sum has only \(i=0\), so only \(\mu=0\)
is allowed. If \(c=0\), the left polynomial vanishes because
\(d-j\ge\lfloor N/2\rfloor>0\).

\begin{proposition}[Exact continuum envelope]\label{lc:left-maximum}
For \(c>0,b\ge0\),
\begin{equation}\label{lc:left-envelope}
 \sup_{\substack{0\le\delta\le1/2\\0\le\mu\le\nu\le1/2}}
 F_{\mathrm L}(\mu,\nu,\delta)
       =\max\{\log c,\mathcal Q(x,b,y)+\tfrac12\log c\}.
\end{equation}
For \(b=0\) restrict the left side to \(\mu=0\).
This is an equality of continuum upper envelopes, not a lower
bound for an oscillatory finite polynomial.
\end{proposition}
\begin{proof}
First assume \(b>0\), \(0<\delta<1/2\), and set
\(z=\sin^2(\pi\delta)\), \(\mathcal D=A^2-4Pz\).
The ratio calculation in Lemma~\ref{lc:cofactor-max} shows that
\(L_\delta\) increases from zero to infinity before \(t_c\) and
decreases from infinity to
\[
                         L_{\rm end}=\frac{A+\sqrt{\mathcal D}}{2z}
\]
afterward. The endpoint derivatives are
\[
 \partial_\mu F_{\mathrm L}=\log b-\log L_\delta(\mu),\qquad
 \partial_\nu F_{\mathrm L}=\log L_\delta(\nu)-\log c.
\]
A stationary \(\nu<t_c\) is a minimum, and a stationary \(\mu>t_c\)
is a minimum. The logarithmic singularity creates no maximum:
the derivative in \(\nu\) is positive on both sides near \(t_c\),
and that in \(\mu\) is negative. The boundary \(\mu=0\) is excluded
for \(b>0,\nu>0\) by its positive initial derivative.
Apart from the diagonal \(\mu=\nu\), the possible maxima therefore
have \(\mu=\mu_b\) and either \(\nu=1/2\) or a stationary
\(\nu_c>t_c\).

The diagonal contributes at most
\(\max(\log c,\frac12\log b+\frac12\log c)\).
The terminal value is \(Q_b(\delta)+\frac12\log c\), because the
complete prefix integral is zero. Both are bounded by the proposed
answer, using Lemma~\ref{lc:action-identity} and
\(\mathcal Q(x,b,y)\ge\frac12\log b\).

A post-change stationary endpoint exists precisely when
\(c\ge L_{\rm end}\). If \(c\le\max(x,y)\), this never occurs
in the open interval. Otherwise it begins at \(\delta_c\), where
\begin{equation}\label{lc:nu-c}
 \sin^2(\pi\delta_c)=\frac Ac-\frac P{c^2},\qquad
 \sin^2(\pi\nu_c)
       =\frac{c^2\cos^2(\pi\delta)}{(c-x)(c-y)}.
\end{equation}
At entry \(\nu_c=1/2\), and afterward \(t_c<\nu_c<1/2\).
Define
\[
 J_c(\delta)=\nu_c\log c+
 \int_0^{\nu_c}\log|a_\delta(t)|\,dt+
 \int_{\nu_c}^{1/2}\log\lambda_\delta(t)\,dt.
\]
The straddling value is exactly
\begin{equation}\label{lc:straddling}
                  S_{\mathrm L}=Q_b(\delta)+\log c-J_c(\delta).
\end{equation}
Subtracting the displayed integrals and using
\eqref{lc:prefix-zero} proves this with no additive constant.

Put \(T_c^-=\sqrt{c^2z-cA+P}\) and
\(\mathcal H_c=2cz-A+2\sqrt z\,T_c^-\). On the strict branch,
\begin{equation}\label{lc:J-derivative}
                  J_c'=-\log(\mathcal H_c/\sqrt{\mathcal D}).
\end{equation}
Here is a normalization proof. The stationary envelope rule gives
\(\partial_{\log c}J_c=\nu_c\), while direct differentiation gives
\[
 \partial_\delta\nu_c=-\frac{c\sin(\pi\delta)}{T_c^-},\qquad
 \partial_{\log c}\log\mathcal H_c=
                         \frac{c\sin(\pi\delta)}{T_c^-}.
\]
Thus the difference in \eqref{lc:J-derivative} is independent of
\(c\) for each \(\delta\). At \(c\downarrow L_{\rm end}\), one has
\(\nu_c\to1/2\), \(T_c^-\to0\), and \(\mathcal H_c\to\sqrt{\mathcal D}\).
The complete prefix integral has zero derivative. The remaining
tail shrinks where \(a_\delta(1/2)\ne0\), so \(J_c'\to0\).
This fixes the constant. To justify differentiation despite the
interior logarithmic zero, write \(\ell(t)=\log|2\cos(\pi t)|\).
At fixed \(\nu>t_c\), the prefix integral is
\[
 \int_\delta^{\nu+\delta}\ell(t)\,dt+
 \int_{-\delta}^{\nu-\delta}\ell(t)\,dt.
\]
Its derivative is
\(\ell(\nu+\delta)-\ell(\nu-\delta)\); the two terms at
\(\pm\delta\) cancel. The moving zero lies strictly inside the
integration interval, while all its endpoints are regular.
Thus this primitive differentiates normally on the strict branch.
At \(\nu=1/2\) its derivative is zero. The tail integral has
smooth integrand there, so the stated branch-entry limit follows.

It follows that \(S_{\mathrm L}'=\log(\mathcal H_c/\mathcal E_b)\).
Both positive quantities are larger roots with the same product
\(\mathcal D\):
\[
\begin{aligned}
 \mathcal E_b&=C_b+\sqrt{C_b^2-\mathcal D},&C_b&=A+2bz,\\
 \mathcal H_c&=C_c+\sqrt{C_c^2-\mathcal D},&C_c&=2cz-A.
\end{aligned}
\]
Both \(C_b,C_c\ge\sqrt{\mathcal D}\). The larger-root function is
increasing, and hence
\[
             \operatorname{sign}S_{\mathrm L}'
                      =\operatorname{sign}((c-b)z-A).
\]
As \(z\) increases, this sign can pass only from negative to positive.
There is no interior spectral maximum. At branch entry the value
is the terminal value already bounded. At \(\delta\to1/2\), both
\(\mu_b,\nu_c\to0\), and logarithmic integrability gives
\(S_{\mathrm L}\to\log c\). The upper bound follows.

The continuum candidates are attained: \(\mu=\nu=0\) gives \(\log c\),
and \(\delta=0,\nu=1/2,\mu=\mu_b\) gives
\(\mathcal Q(x,b,y)+\frac12\log c\).
For \(b=0\), use its sole start or the continuous limiting formulas.
For \(P=0\), the exact one-step root gives the same identities
without division by \(P\). The exceptional \(\mathcal D=0\) occurs
only at \(x=y,\delta=1/2\), where the computed continuous endpoint
values apply.
\end{proof}

